\documentclass[11pt,a4paper]{article}
\usepackage[utf8]{inputenc}
\usepackage[T1]{fontenc}
\usepackage[brazil]{babel}
\usepackage[margin=2.5cm]{geometry}
\usepackage{mathptmx}
\usepackage{enumitem}
\usepackage{booktabs}
\usepackage{tabularx}
\usepackage{array}
\usepackage{hyperref}
\hypersetup{colorlinks=false, pdfborder={0 0 0}}
\usepackage{fancyhdr}

\pagestyle{fancy}
\fancyhf{}
\fancyhead[L]{\small TrackFleet --- Modelo e Escalonamento}
\fancyhead[R]{\small\thepage}
\renewcommand{\headrulewidth}{0.4pt}

\newcolumntype{L}{>{\raggedright\arraybackslash}X}
\newcolumntype{C}{>{\centering\arraybackslash}X}
\newcolumntype{R}{>{\raggedleft\arraybackslash}X}
\setlist{topsep=3pt}

% Cabeçalho padrão do projeto. No documento consolidado (trabalho_pratico_01/)
% estes três comandos são redefinidos e este mesmo arquivo vira parte do capítulo do domínio.
\newcommand{\cabecalho}[3]{%
  \begin{center}
  {\Large\bfseries DOCUMENTO DE #1 DO PROJETO}\\[2pt]
  {\large\bfseries TrackFleet --- Gestão de Rastreador de Automóveis}\\[4pt]
  Disciplina de Gerenciamento de Projetos $\cdot$ Framework PMBOK\textregistered\ 8\textsuperscript{a} Edição $\cdot$ Domínio: #2\\[2pt]
  Integrantes: Enzo Allebrand e Kerlon Ribeiro (dois GPs) $\cdot$ 4\textsuperscript{o} semestre $\cdot$ Data: #3
  \end{center}
  \vspace{2mm}\hrule\vspace{4mm}}
\newcommand{\parte}[1]{\noindent\textbf{\large #1}\par\vspace{1mm}}
\newcommand{\separador}{\par\vspace{2mm}\hrule\vspace{4mm}}

\begin{document}

\cabecalho{GOVERNANÇA}{Governança}{13/08}

\parte{Modelo, Autoridade e Escalonamento}
\noindent Como e por quem as decisões do projeto são tomadas, do dia a dia aos impasses que precisam de uma palavra final.

\section{Definições}
\begin{itemize}[itemsep=4pt]
  \item \textbf{Modelo escolhido:} híbrido e leve; os dois GPs tomam as decisões do dia a dia, e o patrocinador aprova apenas os marcos e as mudanças relevantes.
  \item \textbf{O que a dupla decide sozinha:} ferramentas e tecnologias utilizadas, divisão de tarefas entre os integrantes e decisões técnicas de implementação. Cada GP decide na sua frente, conforme a matriz RACI; o que cruza as duas frentes é decidido por consenso.
  \item \textbf{O que precisa subir ao patrocinador:} mudanças no objetivo do projeto, no prazo de entrega ou no escopo essencial (por exemplo, trocar o público-alvo de transportadoras para uso pessoal), e qualquer limite de risco ultrapassado, conforme o Plano de Gerenciamento dos Riscos.
  \item \textbf{Escalonamento:} sem consenso entre os dois GPs, a decisão sobe direto ao patrocinador, que desempata; se a decisão afetar objetivo, prazo ou valor esperado, também escala para o patrocinador.
  \item \textbf{Cadência de decisão:} check-in semanal de 10 minutos no início de cada encontro da disciplina, para revisar pendências, decisões em aberto, indicadores e riscos. A sincronização diária curta do Termo da Equipe trata só de bloqueios; o que for decidido nela entra no registro de decisões.
\end{itemize}

\end{document}
